\textbf{Neutrino emission of the supernova:}

From the table of constants, the core mass
$M = 1M_\odot = 1.988 \cdot 10^{30}$\,kg.
Mass of one atom of iron $^{56}_{26}\mathrm{Fe}$:
\[
  m_{\mathrm{Fe}} = 56\,\mathrm{Da}
  = 56 \times 1.661 \times 10^{-27}\,\mathrm{kg}
  = 9.3016 \times 10^{-26}\,\mathrm{kg}.
\]
(Note: atomic mass is expressed in \emph{atomic mass units}, symbol a.m.u.\ or
u, also called \emph{Daltons}, symbol Da.)

The number of $^{56}_{26}\mathrm{Fe}$ nuclei in the core is therefore:
\[
  n_{\mathrm{Fe}} = M/m_{\mathrm{Fe}} = 2.1373 \times 10^{55}
\]
\hfill(1 points)

and the initial number of nucleons is:
\[
  n_{\mathrm{nuc}} = 56\,n_{\mathrm{Fe}} = 1.1969 \times 10^{57}.
\]
\hfill(1 points)

Correspondingly, the initial number of protons is
$n_{\mathrm{p}} = 26\,n_{\mathrm{Fe}}$ and the initial number of neutrons is
$n_{\mathrm{n}} = 30\,n_{\mathrm{Fe}}$. \hfill(2 points)

We assume all the nucleons in the original $^{56}_{26}\mathrm{Fe}$ nuclei are
converted to individual nucleons in the given proportions
(1 proton : 8 neutrons). Therefore after the explosion,
\[
  n'_{\mathrm{n}} = \frac{8}{9} n_{\mathrm{nuc}} = 1.0634 \times 10^{57}
\]
neutrons and
\[
  n'_{\mathrm{p}} = \frac{1}{9} n_{\mathrm{nuc}} = 1.3299 \times 10^{56}
\]
protons remain, and thus
\[
  n_{\mathrm{p}} - n'_{\mathrm{p}} = 26\,n_{\mathrm{Fe}} - n'_{\mathrm{p}}
  = 4.2271 \times 10^{56}
\]
\hfill(4 points)

protons are changed. Since 1 neutrino is produced for every converted proton
($\mathrm{p} + \mathrm{e}^- \to \mathrm{n} + \nu_{\mathrm{e}}$), the neutrino
flux is:
\[
  I_\nu = \frac{4.2271 \times 10^{56}}{0.01\,\mathrm{s}}
  = 4.2271 \times 10^{58}\ \mathrm{neutrinos\ s^{-1}}.
\]
\hfill(2 points)

The flux density observed at Earth is given by:
\[
  F_{\mathrm{SN}} = \frac{I_\nu}{4\pi d_{\mathrm{Gal}}^2}.
\]
Substituting
$d_{\mathrm{Gal}} = 8\ \mathrm{kpc} = 8000 \times 3.086 \times 10^{16}\,\mathrm{m}
= 2.4688 \times 10^{20}\,\mathrm{m}$,
\[
  F_{\mathrm{SN}} \approx 5.5 \times 10^{16}\ \mathrm{neutrinos\ s^{-1}\,m^{-2}}.
\]
\hfill(2 points)

\textbf{Neutrino flux from the Sun:}

We can assume that the primary source of Solar luminosity is the p--p reaction,
\[
  4\mathrm{p}^+ + 2\mathrm{e}^- \to {}^4\mathrm{He}^{2+} + 2\nu_{\mathrm{e}}.
\]
Neglecting the electrons and neutrinos, the change in mass $\Delta m$ is:
\[
  \Delta m = 4 \cdot 1.673 \times 10^{-27}\,\mathrm{kg}
  - 6.645 \times 10^{-27}\,\mathrm{kg}
  = 4.7 \times 10^{-29}\,\mathrm{kg},
\]
and thus from $E = (\Delta m)c^2$ the amount of energy released is
${\approx}\,4.2244 \times 10^{-12}$\,J. Solar luminosity is
$L_\odot = 3.826 \times 10^{26}$\,W, so the number of reactions per second is
approximately:
\[
  \frac{3.826 \times 10^{26}}{4.2244 \times 10^{-12}} = 9.057 \times 10^{37}.
\]
With 2 neutrinos released per reaction, this gives a Solar neutrino flux of:
\[
  I_{\nu\odot} \approx 1.8 \times 10^{38}\ \mathrm{neutrinos\ s^{-1}}.
\]
(If the student remembers that the p--p reaction releases
$26.73\ \mathrm{MeV} = 4.2826 \times 10^{-12}$\,J, they will get
$8.9338 \times 10^{37}$ reactions per second and the same final answer.)
\hfill(getting to $I_{\nu\odot}$: 5 points)

At the distance of Earth, the observed flux density is given by:
\[
  F_\odot = \frac{I_{\nu\odot}}{4\pi d_\odot^2},
\]
where $d_\odot = 1\ \mathrm{au} = 1.496 \times 10^{11}$\,m, and so:
\[
  F_\odot \approx 6.4 \times 10^{14}\ \mathrm{neutrinos\ s^{-1}\,m^{-2}}.
\]
\hfill(1 point)

The final ratio is therefore:
\[
  F_{SN}/F_\odot \approx \frac{5.5 \times 10^{16}}{6.4 \times 10^{14}}
  \approx 86 \approx 100,
\]
or two orders of magnitude. \hfill(2 points)
